% ======================================================================
\section{Introduction}\label{sec:intro}

\subsection{Motivation and the gap}\label{sec:need}
More than 430 million people live with disabling hearing loss~\cite{who2021hearing}. Deaf and hard-of-hearing (DHH) viewers receive speech through subtitles. The rest of the soundtrack is mostly lost.
Captions for the deaf may name a sound in brackets, but the DCMP Captioning Key~\cite{dcmp} lets captioners leave out a
sound whose source is clearly seen on screen, and a study of more than 715{,}000 YouTube videos
found hand-written non-speech captions in only 5--7\,\% of them~\cite{may2024unspoken}. DHH users want to know about
the sounds that carry information -- knocks, alarms, sirens, a baby crying~\cite{jain2019home,soundwatch} -- and want
their type and timing~\cite{may2025choices}. Among the needed sounds in this report's benchmark, 30 of 123 (about
a quarter) are danger sounds or key moments (importance 3, Section~\ref{sec:needed}).

What a viewer loses therefore depends on whether the source is visible. A siren over a visible ambulance tells the
viewer nothing new. A siren over a crowd looking up \emph{is} the scene. This project targets the second case, which
none of the existing tools covers. Subtitles carry speech only. Hand-written sound captions are rare; YouTube's automatic captions add only three sound tags ([MUSIC], [APPLAUSE],
[LAUGHTER])~\cite{chaudhuri2017youtube}, whatever is on screen; and the captioning rules
say which sounds to describe but give no tool to find them. Sound-awareness apps for DHH
users~\cite{soundwatch,homesound} alert to sounds in the room, not in a recorded video, and do not ask whether the
viewer can already see the source. Automatic audio captioning~\cite{drossos2017aac,kim2019audiocaps,drossos2020clotho}
describes the whole soundtrack, visible and invisible alike. No system found in the literature reviewed in
Section~\ref{sec:related} shows exactly the sounds the viewer cannot see, at the moment they happen.

\subsection{Problem}\label{sec:problem}
Given a video, produce the same video with a side panel that shows a picture of a sound \emph{while} the sound is
heard and its source is \emph{not} on screen, and shows nothing otherwise. Speech and music are never drawn (speech has
subtitles, and music is marked in ordinary captions); human sounds such as laughter, coughing or a scream are drawn.
The output is a picture rather than a word in brackets because a picture can be taken in at a glance, without
reading, while the viewer is already reading subtitles; this is the design hypothesis behind RQ3
(Section~\ref{sec:rqs}), and it is not tested with viewers here. No model may be trained: the 158 labelled clips are
far too few to train on, and all of them are needed for evaluation. The system must work on one video at a time, as a
captioning tool does when a video is uploaded. For every sound this means four questions:
\begin{enumerate}
  \item \textbf{Is there a sound, and when does it start?}
  \item \textbf{What is it?} At a level a picture can show: \emph{dog bark}, not \emph{animal}.
  \item \textbf{Can the viewer already see it?} If the thing making the sound is visibly making it, stay silent.
  \item \textbf{What should the picture show?}
\end{enumerate}

\subsection{Why ready-made models are not enough}\label{sec:notplug}
At first sight the task looks like a short chain of ready-made models: a sound tagger names the sound, a vision model
says whether its source is on screen, an image model draws it. In practice every link of that chain fails in its own
way, and the failures add up:
\begin{itemize}
  \item \textbf{Hearing.} Taggers rank \emph{Speech} or \emph{Music} above a quiet door or hammer, so the sound never
        becomes a candidate; and of the sounds they do report, only about 40\,\% can be confirmed as real.
  \item \textbf{Confirming.} A sound that only one audio-language model names in its list of heard sounds is real only
        about a third of the time.
  \item \textbf{Seeing.} A vision-language model answers ``is there a bell?'' when asked ``is the bell ringing?'': it takes
        presence for production. Every extra vision tool tried (object boxes, crops, segmentation, synchrony, sound
        localisation) traded one error for another.
  \item \textbf{Drawing.} Audio-to-image models draw the whole soundtrack as one scene and cannot leave out the dog that is
        already on screen.
  \item \textbf{Scoring.} Model judges and model-written references either reward the system that wrote them or decide
        visibility at chance, so even measuring progress needs human per-sound labels.
\end{itemize}
The end-to-end effect is measurable. Even with the full detection stack of the final system, drawing every detected
sound without reading the video (the proposal's direct audio-to-image baseline, Section~\ref{sec:baselines}) finds 26 needed
sounds but shows 61 wrong pictures on the test set, and is no better than showing nothing. Turning the chain into a
system that costs less than every baseline took the rules of Section~\ref{sec:system}, each one answering a failure described in
Section~\ref{sec:hard}: two detectors, one of three independent witnesses for every weak sound, a three-question
on-screen vote, and guarded picture prompts with a check-and-redraw loop. These rules were found after more than 150
scored full-pipeline variants, each judged by the per-sound metric.

\subsection{Research questions}\label{sec:rqs}
From the project proposal (6 August 2026); Section~\ref{sec:conclusion} answers each.
\begin{enumerate}[label=RQ\arabic*.]
  \item Which types of audio information contribute most to scene understanding when presented visually?
  \item What level of semantic granularity should be used to represent ambient sounds?
  \item Can automatically generated visual augmentations improve accessibility beyond subtitles?
  \item How should semantic audio augmentation systems be evaluated?
  \item How does the proposed pipeline compare with representative audio-to-visual approaches?
\end{enumerate}

\subsection{Contributions}
\begin{enumerate}
  \item \textbf{The task, a metric and a benchmark.} A precise rule for when a picture is needed, a per-sound score
        that prices a missed sound against a wrong picture, and a benchmark of 308 timed sounds in 158 video clips,
        each sound marked by whether its source is seen making it, whether it is obvious, and how important it is (Sections~\ref{sec:task}
        and~\ref{sec:benchmark}).
  \item \textbf{A system from open models, without training, that costs less than every baseline of the proposal on both sets.} On
        the test set it costs less than both direct audio-to-image generation ($p < 0.001$) and showing nothing
        ($p = 0.034$, 0.068 after Holm): clearly against direct audio-to-image generation, at the margin against
        showing nothing. Audio captioning scores the same as direct audio-to-image generation under the per-sound metric. On the
        test set the system keeps every danger sound that drawing every sound finds (10 of 13) while showing 38 fewer wrong pictures
        (Sections~\ref{sec:system} and~\ref{sec:results}).
  \item \textbf{Two findings on why ready-made models are not enough.} Vision-language models take presence for
        production, and every extra vision tool trades one error for another; and no audio stage tried here separates
        a right off-screen name from a wrong one reliably: most rules that added hits also added two or more wrong pictures per new hit
        (Section~\ref{sec:hard}, with the capabilities today's models lack in Section~\ref{sec:missing}).
  \item \textbf{A pre-registered method at scale.} 66 pre-registered rounds, each accepting a change only by a rule
        written down before its run, over more than 150 full-pipeline variants and more than 30 open models
        (Section~\ref{sec:scale}, Appendix~\ref{app:history}), with a record of every round (Table~\ref{tab:rounds}) and the main rejected ideas
        (Table~\ref{tab:rejected}), so the next attempt need not repeat them. No single model swap produced the gain;
        it is the sum of small rules, each closing one failure.
  \item \textbf{An evaluation protocol} (the answer to RQ4): per-sound human labels, errors split by cause, and a cost
        reported over a range of weights (Sections~\ref{sec:task}, \ref{sec:sensitivity} and~\ref{sec:analysis});
        model-written references and model judges were tried and fail at the visibility question
        (Section~\ref{sec:hard-measure}).
  \item \textbf{A released, exactly reproducible system and record.} Code, per-sound labels, scoring harness and the
        logged decision of every candidate on all 158 clips (release v1.3.10); every count of the final system reproduced by two
        independent routes with no difference on any clip (Section~\ref{sec:support}); the source file of every number
        (Table~\ref{tab:sources}); and the full pre-registration record in release v1.2.0, so each claim can be checked at
        the level of a single sound (Section~\ref{sec:engineering}, Appendix~\ref{app:sources}).
\end{enumerate}

\paragraph{How the report is organised.} Section~\ref{sec:related} reviews related work.
Sections~\ref{sec:task}--\ref{sec:system} define the task, the data and the system; Section~\ref{sec:results} gives the
results, Section~\ref{sec:hard} explains why the problem is hard for today's models and Section~\ref{sec:analysis}
gives the evidence on the remaining errors; Section~\ref{sec:limits} lists the limits and threats to validity, and Section~\ref{sec:conclusion} answers the research questions. The
appendices hold the prompts, the detection details and configuration, reproduction steps with the source of every
number, a glossary of internal names, and the full development history. The code, labels and report sources are
published as release v1.3.10 of the project repository (see
\hyperref[sec:availability]{Code and data availability}).
